\documentclass[a4paper,12pt,titlepage]{article}
\usepackage{graphicx}
\usepackage[hidelinks]{hyperref}
\usepackage{listings}
%\setcounter{tocdepth}{1}
\usepackage{float}




\DeclareGraphicsExtensions{.png}
%\graphicspath{ {./images/} }

\begin{document}
\begin{titlepage}
	\begin{center}
		
		\begin{figure}[t]
			\centering
			\includegraphics[width=350px]{UP_Logo.png}
		\end{figure}		
		
		\textbf{\LARGE COS301 Mini Project Functional \\Architecture Requirements\\}
		
		\vspace{1 cm}
		
		\LARGE{\textbf{Group Name: }Group 7\_a}
		
		%\begin{minipage}{0.4\textwidth}

		\begin{flushright} \large
			Roger Tavares \emph10167324\newline
			Thinus Naude \emph13019602 \newline
			Kabelo Kgwete \emph11247143\newline
			Sylvester Mpanganer \emph11241617\newline
			Maphuti Setati \emph12310043\newline
			Ruth Ojo \emph12042804\newline
			Axel Ind \emph12063178\newline
			Lindelo Mapumulo \emph12002862\newline
			Maria Qumalo \emph29461775\newline
		\end{flushright}
		%\end{minipage}
		
		\vfill
		
		\textbf{Git repository link:\\}
		\url{https://github.com/thinusn/COS301MiniProjectArchitectureRequirements}
		
		\vfill
		
		{\LARGE Final Version}
		\\
		{\large \today}		
		
		
	\end{center}
\end{titlepage}



%-----------------TABLE OF CONTENT-----------------
\newpage
\tableofcontents


%-----------------BASIC INTRODUCTION-----------------
\newpage
\section{Introduction}

This document was compiled by our group during our meetings and was produced as a whole by the team.\\\\

This document contains specifications of the software architecture 
\\requirements. This is the infrastructure upon which the application
\\ functionality will be developed. The following non-functional
\\ requirements are addressed in depth with supporting diagrams
\\(when necessary):

\begin{itemize}
	\item Quality requirements.
	\item Integration and access channel requirements.
	\item Architectural constraints
	\item Architectural patterns or styles
	\item Architectural tactics or strategies
	\item Use of reference architectures and frameworks
	\item Access and integration channels
	\item Technologies
\end{itemize}

%-----------------PARTS TO EDIT-----------------
\newpage

\section{Architecture Requirements}
\subsection{Architectural Scope}
\newpage
\subsection{Critical Quality Requirements}
\subsubsection{Scalability}%Ruth
	\subsubsection*{Description}
	Scalability is an essential aspect of a system and is the ability of a system to be easily enlarged in order to accommodate a growing amount of work.
	\subsubsection*{Justification}
	 As the Buzz System should allow for more than a million concurrent users, the system must be able to handle that.
	\subsubsection*{Mechanism}
		\begin{enumerate}
			\item Strategy:\\\\
			Scalability can be achieved by:
			\begin{itemize}
			\item Clustering: using more resources by running many instances of the application over a cluster of servers or instances.
			\item Load Balancing: Spread the systems load across time or across resources through scheduling or queueing. In the case that the limit for a server has been reached, a new instance will have to be created in order to handle the number of increasing requests. On the other hand, if the usage of a server is way below the capacity, the number of instances will have to be reduced.\\\\
			In terms of databases:
			\item Caching: to ensure no duplication or repeated retrieval of frequent objects or queries; a separate module can facilitate caching; thus system resources will not be used up unnecessarily.
			 \end{itemize}
			\item Architectural Pattern(s):
			\begin{itemize}
			\item Blackboard Architectural Pattern
				\begin{itemize}
					\item Makes used of a spaced-based architecture, making it easy to add a new knowledge source or update an existing knowledge source.
				 \end{itemize}
			 \end{itemize}
		\end{enumerate}

\newpage
\subsubsection{Security}%Kabelo
<<<<<<< HEAD
	Description: \\\\
	Justification: \\\\
	Mechanism:
	\begin{enumerate}
		\item Strategy: 
		\item Architectural Pattern:
	\end{enumerate}
\newpage
=======
	\subsubsection*{Description}
	Security includes the securing of confidential data stored from unauthorized access and modification. For BuzzSpace, this means:
	\begin{itemize}
		\item Ensuring only genuine users may have access to the system’s data. 
		\item Different levels of access should be provided to users.
		\item Securing vulnerabilities that may exist in the application.
	\end{itemize}	
	\subsubsection*{Justification}
	 Security is a crucial aspect of the system. Poor security design and implementation can compromise the data stored by the system and hackers with malicious intentions can succeed to gain access to sensitive information. Using a combination of security strategies and patterns, we can prevent unauthorized users from accessing sensitive information or cause damage to the system.
	\subsubsection*{Mechanism}
		\begin{enumerate}
			\item Strategy:\\\\
			Security can be achieved by using:
			\begin{itemize}
			\item Authentication: This is the process that ensures and confirms a user’s identity.
			\item Encryption: This process converts sensitive data into another form, which cannot be easily understood by unauthorized users. Data will be converted to readable format for authorized users.
			\item Input Validation: The process of validating all the input to an application before using it. This will prevent malicious attacks such as SQL injections.
			\item Auditing and Logging: This includes the recording of the implementation level events and domain-level events that occur when the system is being used. This data will help identify what went wrong if the system experiences malicious attacks.\\\\
			
			 \end{itemize}
			\item Architectural Pattern(s):
			\begin{itemize}
			\item Authentication Enforcer pattern
				\begin{itemize}
					\item This pattern can be used to manage and delegate authentication processes.
				 \end{itemize}
			\item Secure Base Action pattern
				\begin{itemize}
					\item Shows centralizing handling of security tasks such as input validation in a base action class.
				 \end{itemize}
			\item Secure Logger pattern
				\begin{itemize}
					\item This pattern can be used to log sensitive data and ensuring tamper-proof storage.
				 \end{itemize}
			\item Secure Session Manager
				\begin{itemize}
					\item This pattern can be used to show securely centralizing session information handling.
				 \end{itemize}
			 \end{itemize}
		\end{enumerate}

>>>>>>> origin/master
\subsubsection{Usability}%Lindelo
	Description: Usability is critical because Buzz is user-oriented and requires that users don't get frustrated with the interface.\\\\
	Justification: Buzz should be easy to use and the user should be able to remember how Buzz works; despite their long absence. The interface should appear easy to use; the visuals should not frustrate the user.\\\\
	Mechanism:
	\begin{enumerate}
		\item Strategy:
		\begin{itemize}
			\item Descriptive headings that will not be ambiguous to the user.
			\item Clear and descriptive "help" messages that will initiate the user into Buzz. 
			\item Descriptive error messages that will tell the user an error has occurred and the necessary steps to address the errors.
			\item Grouped content; this will make the interface look easier to use. If a user is looking for content, they should find related content in the same place. Navigation becomes easier.
			\item A user should be given feedback on their actions. If a page is still loading, for example, the user should be aware that the page is loading.
		\end{itemize} 
		\item Architectural Pattern:
	\end{enumerate}'
\newpage
\subsubsection{Integrability}%Sylvester
	
\subsubsection*{Description:}
	  The Buzz system should be integrated with ease on any host site.
	 \subsubsection*{Justification:}
	  The Buzz system is primarily designed for the CS website, but if in future a need to host it on another site arise,it needs to be easily integrated to the the site without any extra cost and the need  to change the implementation of the system.
	\subsection*{Mechanism:}
	\begin{enumerate}
		\item Strategy: \\\ The contracts based decoupling strategy can be used to implement integrability and achieve extensibility on the system by providing an interface between the system and the host site which will improves decoupling on the host site and the system.    
		\item Architectural Pattern: \\\ Messaging can be used to enhance decoupling and increase flexibility, the Buzz system or host site can change without affecting the other party, both parties will act as independent component that use an interface between them to request services and share resources.  
	\end{enumerate}

\newpage
\subsection{Important Quality Requirements}

\subsubsection{Performance}%Thinus
	\subsubsection*{Description}
		Performance is an indication of the responsiveness of a system to execute some action(s) within a specific time interval. This time interval is usually measured in terms of latency or throughput.
		\begin{quote}
			\textbf{Latency:} the time it takes the system to respond to any event(s).
		\end{quote} 
		\begin{quote}
			\textbf{Throughput:} the number of events executed within a given amount of time by the system.
		\end{quote} 
		A system with high performance maximizes throughput and minimizes latency. This is ideally what we want for the BuzzSystem
	\subsubsection*{Justification}
	 	We chose performance as an important quality requirement because the over/under performance of a system will influence other quality attributes of the system. If the system has increased latency and decreased throughput( for example due to the systems inability to handle increased load, or subpar scalability) the performance, responsiveness and usability of the system will suffer. Because the Buzz System will be used by users we want them to use a system the performs optimally and responds to their event(s) with minimal latency. 
	\subsubsection*{Mechanism}
		\begin{enumerate}
			\item Strategy:
			The system needs to be responsive. The time it takes for a request to be sent back to a user should be minimized.
			\begin{itemize}
				\item The system needs to be hardware and software fault tolerant (to a certain degree). The system needs to continue working/running as intended but possibly at a reduced level, rather than breaking/stopping completely. A system that is off-line has 0 throughput. 
				\item The UI presented to the user needs to be clean, dynamic and use minimal resources. For example using minified JavaScript files and compressed images. In order to reduce latency.
				\item The performance of the system needs to be optimal, memory or processor intensive tasks should run/execute when there are the least number of active users in order to minimize the impact these tasks will have on performance. For example run batch plagiarism and netetiquette checks at 2AM in the morning.
				\item The system needs to have a coping mechanism when there is a sudden change in its environment. For example if the system can currently handle 100 connections and suddenly there is 10000 connections. the system needs to know how to cope. Performance will suffer if this is not taken into consideration. See Scalability for possible solution.
				\item The system should deliver intermediate results or updates to the user when executing a request. For example, a web page that submits a form via AJAX can have a status/busy indicator to let the user know his/her request is being processed.
				\item When there is increased database server processing the system needs to keep latency low and throughput high. The system can achieve this by storing frequently requested database results or objects in a cache. This will ensure that frequent objects and queries aren't repeated or fetched over and over, wasting system resources and increasing database server processing.
				\item An important aspect of performance is to ensure the scalability of the system is optimal. Because an increase in performance directly affects the systems scalability, and in turn the lack of scalability also affects performance.
			 \end{itemize}
			\item Architectural Pattern(s):
			\begin{itemize}
			\item Many/all of the patterns discussed in this document could assist in performance.
			\item Architectural Patterns that allow for scaling of resources are preferred. 
			\end{itemize}
		\end{enumerate}

	\newpage
\subsubsection{Plug-ability (Maintainability)}%Roger
	\input{./plug-ability.tex}
\newpage
\subsubsection{Monitor-ability}%Maria
	Description: \\\\
	Justification: \\\\
	Mechanism:
	\begin{enumerate}
		\item Strategy: 
		\item Architectural Pattern:
	\end{enumerate}
\newpage
\subsection{Nice to have quality requirements}
\subsubsection{Reliability and Availability}%Axel
	Description: \\\\
	Justification: \\\\
	Mechanism:
	\begin{enumerate}
		\item Strategy: 
		\item Architectural Pattern:
	\end{enumerate}
\subsubsection{Testability}%Maphuti
	Description: \\\\
	Justification: \\\\
	Mechanism:
	\begin{enumerate}
		\item Strategy: 
		\item Architectural Pattern:
	\end{enumerate}
\newpage
\subsection{Architectural Constraints}
	\begin{itemize}
		\item\textbf{Neo4j}, the preferable database technology, is schema-less, which means a schema change will not affect how the data is stored. It includes other features that would improve Buzz, but JQPL has to be used. \\ 
		\item\textbf{JPQL} was chosen by the client, although Neo4j would have been better. \\ 
		\item\textbf{Java EE} is the platform that will be used to implement Buzz, the system is constrained to this platform. \\ 
		\item\textbf{JSF}, was chosen by the client and forms part of Java EE. \\ 
		\item\textbf{AJAX} the client requested that this technology be used.  \\ 
		\item\textbf{HTML} \\
		\item\textbf{The operating system} that will run the Buzz server is Linux, as requested by the client.
		\end{itemize}
\newpage
\section{Architectural Patterns or Styles}%Axel
\newpage
\section{Architectural Tactics or Strategies}%Kabelo
\newpage
\section{Use of Reference Architectures and Frameworks}%Thinus
\newpage
\section{Access and Integration Channels}
\subsection{Access Channels}
\subsubsection{Human Access Channels}%Maria
\subsubsection{System Access Channels}%Roger
	System Access Channels is all the external services or applications that have to interact with the Buzz system in some way or another to ensure that all functionality is available to the users. The main systems that will be used in the current project scope are:

\begin{enumerate}
	\item The UP Computer Science Website:
	\begin{itemize}
		\item The Computer Science website must be able to integrate with the system allowing students to access it from module pages providing these users with help and support in their studies.
		\item This system will eventually replace the current outdated CS forum that is being used.
	\end{itemize}
	\item LDAP:
	\begin{itemize}
		\item LDAP is the database used by the University to store and record all of the students information. This database can be modified or expanded to hold user ranks and post information.
		\item This means that access and permissions to the database by the Buzz system need to be provided by the database administrators. 
	\end{itemize}
	\item The Internet:
	\begin{itemize}
		\item Since this system will be integrated with the CS website as mentioned above it will in turn also require access to the Internet allowing students to not only access the system locally but also at home or remotely as mentioned in Human Access Channels.
	\end{itemize}
	\item External Plagiarism Checker:
	\begin{itemize}
		\item Developing a plagiarism checker from scratch is time consuming and also outside the scope of the system so in order to assess the posts for potential copyright infringement access to an external checker will need to be incorporated.
		\item TurnItIn is a company that provides such a service.
	\end{itemize}
\end{enumerate} 
\subsection{Integration Channels}%Sylvester
	\subsubsection*{Channels}
\begin{itemize}
\item The Buzz system will access the CS MySQL database to retrieve information of a particular module/course.
\item The Buzz system will access the CS LDAP server to gain access to and retrieve details of students,  lectures and also a module's enrollment  list.
\end{itemize}

\subsubsection*{Protocols}
The Buzz system will mainly use the HTTPS, LDAP and SMTP protocols.
\begin{itemize}
\item HTTPS: This protocol will be use for security, to ensure that a secure connection is maintained and users information is safe. 
\item LDAP: The protocol that will be used for the integration of the system with the CS LDAP server.
\item SMTP: Notifications for plagiarism or posts not conforming to the netiquette rules will use this protocol for the Buzz email system to enable the communication between the administrator (lecture) and students.  
\end{itemize}
\newpage
\section{Technologies}%Lindelo and Ruth


\end{document}
